\documentclass{article}

\usepackage{fancyhdr}
\usepackage{extramarks}
\usepackage{amsmath}
\usepackage{amsthm}
\usepackage{amsfonts}
\usepackage{tikz}
\usepackage[plain]{algorithm}
\usepackage{algpseudocode}
\usepackage{color}

\usetikzlibrary{automata,positioning}

%
% Basic Document Settings
%
%%% page layout
\topmargin=-0.45in
\evensidemargin=0in
\oddsidemargin=0in
\textwidth=6.5in
\textheight=9.0in
\headsep=0.25in

\linespread{1.1}    %%% line spacing

\definecolor{ustcblue}{cmyk}{1,0.8,0,0}

\pagestyle{fancy}
\lhead{\hmwkAuthorName}
\chead{\hmwkClass\ (\hmwkClassInstructor): \hmwkTitle}
\rhead{}
\lfoot{\lastxmark}
\cfoot{\thepage}

\renewcommand\headrulewidth{0.4pt}
\renewcommand\footrulewidth{0.4pt}

\setlength\parindent{0pt}

%
% Create Problem Sections
%

\newcommand{\enterProblemHeader}[1]{
    \nobreak\extramarks{}{Problem \arabic{#1} continued on next page\ldots}\nobreak{}
    \nobreak\extramarks{Problem \arabic{#1} (continued)}{Problem \arabic{#1} continued on next page\ldots}\nobreak{}
}

\newcommand{\exitProblemHeader}[1]{
    \nobreak\extramarks{Problem \arabic{#1} (continued)}{Problem \arabic{#1} continued on next page\ldots}\nobreak{}
    \stepcounter{#1}
    \nobreak\extramarks{Problem \arabic{#1}}{}\nobreak{}
}

\setcounter{secnumdepth}{0}
\newcounter{partCounter}
\newcounter{homeworkProblemCounter}
\setcounter{homeworkProblemCounter}{1}
\nobreak\extramarks{Problem \arabic{homeworkProblemCounter}}{}\nobreak{}

%
% Homework Problem Environment
%
% This environment takes an optional argument. When given, it will adjust the
% problem counter. This is useful for when the problems given for your
% assignment aren't sequential. See the last 3 problems of this template for an
% example.
%
\newenvironment{homeworkProblem}[1][-1]{
    \ifnum#1>0
        \setcounter{homeworkProblemCounter}{#1}
    \fi
    \section{Problem \arabic{homeworkProblemCounter}}
    \setcounter{partCounter}{1}
    \enterProblemHeader{homeworkProblemCounter}
}{
    \exitProblemHeader{homeworkProblemCounter}
}

%
% Homework Details
%   - Title
%   - Due date
%   - Class
%   - Section/Time
%   - Instructor
%   - Author
%

\newcommand{\hmwkTitle}{Homework\ \#1}
\newcommand{\hmwkDueDate}{\today}
\newcommand{\hmwkClass}{Inversion}
\newcommand{\hmwkClassInstructor}{Professor H. Yao \& H. Zhang}
\newcommand{\hmwkAuthorName}{\textbf{Jintao Li}}
\newcommand{\hmwkAuthorID}{\textbf{SA20007037}}

%
% Title Page
%

\title{
    \vspace{2in}
    \textbf{\hmwkClass:\ \hmwkTitle}\\
    \normalsize\vspace{0.2in}\large{\hmwkDueDate}\\
    \vspace{0.2in}\large{\textit{\hmwkClassInstructor}}
    \vspace{3in}
}

\author{\hmwkAuthorName \\
\hmwkAuthorID}
\date{}

\renewcommand{\part}[1]{\textbf{\large Part \Alph{partCounter}}\stepcounter{partCounter}\\}

%
% Various Helper Commands
%

% Useful for algorithms
\newcommand{\alg}[1]{\textsc{\bfseries \footnotesize #1}}

% For derivatives
\newcommand{\deriv}[1]{\frac{\mathrm{d}}{\mathrm{d}x} (#1)}

% For partial derivatives
\newcommand{\pderiv}[2]{\frac{\partial}{\partial #1} (#2)}

% Integral dx
\newcommand{\dx}{\mathrm{d}x}

% Alias for the Solution section header
\newcommand{\solution}{\textbf{\large Solution}}

% Probability commands: Expectation, Variance, Covariance, Bias
\newcommand{\E}{\mathrm{E}}
\newcommand{\Var}{\mathrm{Var}}
\newcommand{\Cov}{\mathrm{Cov}}
\newcommand{\Bias}{\mathrm{Bias}}

\begin{document}

\begin{titlepage}

\begin{center}

\textcolor{ustcblue}{\includegraphics[width=0.15\textwidth]{./ustc_logo_fig.pdf} \\ [1cm]}
% Title
{ \huge \bfseries \hmwkClass\ \hmwkTitle}\\[0.4cm]

\hmwkClassInstructor

\hmwkAuthorName \\
\hmwkAuthorID
\vfill
% Bottom of the page
{\large \today}

\end{center}

\end{titlepage}


\begin{homeworkProblem}

Let 
\begin{equation}
\mathbf{A} = 
\left[
\begin{matrix}
  1 & 2 & 3 & 4 \\
  2 & 2 & 1 & 3 \\
  4 & 6 & 7 & 1
\end{matrix}
\right]
\end{equation}
Find bases for $\mathcal{N}(\mathbf{A})$, $\mathcal{R}(\mathbf{A})$, $\mathcal{N}(\mathbf{A}^T)$, 
and $\mathcal{R}(\mathbf{A}^T)$. What are the dimensions of the four subspaces?

\end{homeworkProblem}

\begin{homeworkProblem}

Show that if $\mathbf{x} \perp \mathbf{y}$, then 
\begin{equation}
\left\| \mathbf{x} + \mathbf{y} \right\|^2_2 = 
\left\| \mathbf{x} \right\|^2_2 +
\left\| \mathbf{y} \right\|^2_2 .
\end{equation}

\end{homeworkProblem}

\begin{homeworkProblem}
In this exercise, we will derive the formula (A.88) for the 1-norm of a matrix.
Begin with the optimization problem 
\begin{equation}
\left\| \mathbf{A} \right\|_1 = 
\max_{\| \mathbf{x} \|_1 = 1} \| \mathbf{Ax}\|_1 .
\end{equation}

\part

Show that if $\|\mathbf{x}\|_1 = 1$, then
\begin{equation}
\|\mathbf{Ax}\|_1 \leq 
\max_{j} \sum\limits_{i=1}^m \left| \mathbf{A}_{i, j} \right| .
\end{equation}


\part

Find a vector $\mathbf{x}$ such that $\|\mathbf{x}\|_1 = 1$, and 
\begin{equation}
\|\mathbf{Ax}\|_1 = \max_{j} \sum\limits_{i=1}^m \left| \mathbf{A}_{i, j} \right| .
\end{equation}


\part
Conclude that 
\begin{equation}
\|\mathbf{A}\|_1 = \max_{\|\mathbf{x}\|_1 = 1} \|\mathbf{Ax}\|_1 = 
\max_{j} \sum\limits_{i=1}^m \left| \mathbf{A}_{i, j} \right| .
\end{equation}

\end{homeworkProblem}

\end{document}
